\section{总结与展望}

\subsection{工作总结}

本文围绕“基于投机采样的大模型云端协同技术”这一课题，设计并实现了一个面向 Android 移动端与桌面端协同场景的投机解码实验系统。系统采用“移动端草稿生成、桌面端目标验证”的协同生成模式，Android 端承担草稿提案任务，桌面端承担目标验证任务。

在方案设计方面，本文给出了系统总体结构、投机解码执行流程以及协议与会话管理机制，明确了 \texttt{start}、\texttt{propose}、\texttt{fallback} 和 \texttt{close} 等消息语义，并围绕已接受前缀组织多轮投机会话。系统在实现上完成了草稿会话与目标会话两类状态边界设计，支持统一真实词元空间下的线性提案、树形提案和目标验证。

在实验分析方面，本文以 Android 端侧本地目标模型推理、云侧本地单进程投机加速和云侧本地双进程投机加速作为对照，评估 Android--桌面端跨设备投机解码方案的性能表现。实验结果显示，跨设备方案在长输出任务中完成了多轮投机会话，吞吐率明显高于端侧直接运行目标模型，并取得了接近云侧本地双进程方案的表现。进一步分析表明，跨设备方案与云侧本地方案的差异主要来自 Android 草稿侧状态维护、平均提案长度以及跨设备协同边界上的固定开销。

从资源分布看，跨设备方案将草稿模型运行时和部分草稿侧状态迁移到 Android 端，使桌面端主要承担目标验证任务。实验中的内存占用观察显示，该部署方式能够降低桌面端同时承载草稿与目标运行时的资源压力。这一结果说明，端云协同投机解码除生成速度外，还可以从模型运行时分布角度改善云侧资源组织方式。

\subsection{不足与改进方向}

本文完成了系统设计与实验验证，但仍存在以下不足。

首先，统一真实词元表示下的验证机制仍受当前实现条件限制，尚未达到严格输出分布保持的程度。目标侧概率读取、残差修正以及后续词元处理仍受观测窗口限制，与更严格的理论方法相比仍存在差距。

其次，目标侧验证器已具备会话边界和一定持续性，但持久化目标运行时仍需改进。部分验证过程仍需额外的准备或恢复步骤，这会放大端云协同中的中间成本。

再次，草稿侧与目标侧的状态同步仍然限制跨设备投机解码性能。本文采用已接受前缀驱动的状态推进与回滚机制，但在更长上下文、更复杂草稿树或更大模型配置下，这一代价可能上升。

最后，本文实验以系统验证为主。结论的适用范围还需要通过更多模型组合、提示任务和网络条件下的实验加以检验。

\subsection{未来展望}

后续研究可从以下几个方面展开。

\begin{enumerate}[label=\arabic*.]
  \item 完善真实词元空间下的验证逻辑，使目标概率、草稿概率、残差修正和后续词元处理接近严格的输出分布保持要求。
  \item 继续强化草稿运行时与目标运行时的持续化能力，减少每轮投机会话的重复准备和恢复成本。
  \item 研究树形提案与分支条件概率维护机制，提高候选覆盖能力和接受率。
  \item 研究多移动端共享同一云侧目标验证服务的调度机制，使多个端侧设备分别承担草稿生成，并复用云侧常驻目标模型运行时，从而提高云侧模型资源利用率。
  \item 扩展实验维度，包括不同草稿模型与目标模型组合、不同网络环境以及更长上下文任务，评估端云协同投机解码的适用范围。
\end{enumerate}

本文实现的系统完成了移动端草稿生成、桌面端目标验证和多轮投机会话管理。后续工作需要继续围绕词元对齐、运行时连续性、验证精度和多端资源调度降低跨设备协同成本。
